\chapter{Controls, Exploration, and Progression}
\label{written-exploration}

A visit to town gives the party a chance to recover, hear local news, and
prepare for the next trip. Out in the countryside, the same small decisions
continue to matter. Healing before a difficult stretch, keeping useful
items with the right companion, and remembering a convenient resting place
can make a long journey considerably more comfortable.

\subsection*{Game-long achievement plan}

The achievement route has \textbf{70 goals worth 410 points}.
Keep ordinary in-game save slots for the marriage branches and the
separate final-boss lineups. Confirm an unlock in the achievement list
before reloading a protected checkpoint.

\begin{reftable}{P{0.29\linewidth}P{0.07\linewidth}P{0.51\linewidth}}
\toprule
\GuideTableHeading{Achievement} & \GuideTableHeading{Points}
  & \GuideTableHeading{Start early and finish at} \\
\midrule
\endhead
I Don't Want Your Help & 10 & From the start of adulthood, defer the
  wagon and all wild recruitment until Kitty rejoins in 6.2.
  Continue with Hero and Henry, then prepare the Hero for solo travel. \\
Monster Tamer / Recruiter / Hunter / Collector & 2 / 3 / 5 / 10 & After
  Kitty, track distinct recruitable species toward the 1, 5, 10, and
  16 milestones. Confirm each distinct-species milestone in the achievement list. \\
Small Medal Hunter & 25 & Record the 29 fixed searches separately
  from chests and enemy drops. The last search is Shahanna's outer-wall barrel. \\
Small Medal Collector & 25 & Record the 19 fixed chest pickups.
  The last is the Bonus Dungeon entrance chest beside the Restoration Orb. \\
King Metal Shield & 4 & Reserve 43 unspent medals and buy the shield
  at the Medal King. Collected-location progress remains separate from balance. \\
King Metal Sword & 5 & Build 50,000 casino coins through ordinary
  play, then buy the sword. The balance alone is preparation. \\
Let's Cuddle / Special Delivery & 3 / 3 & Preserve the pre-bride
  church save. Finish Bianca's two-person Alcapa stay and Flora's
  two courier gifts on separate branches before Granvania. \\
Kandar's Henchmen & 5 & Defeat the Reinhart and Mountain Village
  guards during reunion visits, then the Shahanna exterior guard. \\
The Sacred Raiment & 5 & Keep Sun Crown, Dragon Staff, and Royal
  Cloak until the Bonus Shield of Radiance. Equip all four on the
  Hero and remove accessories. \\
There and Back Again & 2 & After learning Return, track the actual
  17 destination entries. Finish registration on reaching Evil Mountain. \\
That's Got To Hurt & 5 & Develop a strong focused attacker and
  normal attack support. A single hit must show at least 500 damage. \\
A Family Matter / Heroic Trainer & 10 / 25 & Develop family and
  monster recovery teams separately. Keep the prepared pre-Mildrath
  church save for one restricted lineup through both forms per attempt. \\
Underworld Challenge & 25 & Save the Bonus equipment and develop
  Drum, Barrier, and strong recovery. Defeat Estark within 25 turns
  using ordinary play. \\
\bottomrule
\end{reftable}

\begin{achievementsetup}{Early missable checkpoints}
Defeat Lenoire's Walking Statue before Boss Ghost. Search the slavery
camp's sleeping-room pot before the escape. Complete the nighttime
Reinhart inn inquiry, Ludman's alternative proposal, both marriage
branches, the Hero's failed Helmet fitting, and Flora's courier gifts
before advancing their associated story windows.

Before defeating Ivol, attempt the Staff of Stoross on the wife's statue
and use the purchased Ivol's Book during his battle. Keep the six-item
critical/recruitment shortcut and early-patch trap repetition outside
the ordinary achievement route. The 500-damage and 25-turn goals
explicitly require a no-cheating attempt.
\end{achievementsetup}

\section{Controller inputs, field commands, and menu navigation}

\subsection*{Moving and choosing an action}

Use the directional pad to walk and to move the cursor through menus.
Approach a person or object and face it before choosing an interaction.
A staircase, entrance, or passage generally connects to the next area when
the party walks onto the appropriate square.

Press \textbf{A} to open the field command menu. Choose a command with the
directional pad and confirm it with A. In menus, \textbf{B} cancels the
current selection or closes the current window. Some actions lead through
several choices, such as selecting a companion, an item, and its recipient.
Read the current prompt before confirming the next choice.

\subsection*{The field commands}

\begin{description}
  \item[Talk:] Speak to the person the leading character is facing. Shop
    counters and service desks are approached in the same way.
  \item[Spell:] Choose a companion and one of that companion's field spells.
    Healing, status recovery, and travel spells appear as they become
    available. Select the recipient when the spell requests a target.
  \item[Status:] Inspect a companion's level, attributes, HP, MP, equipment,
    and other personal information. Use this screen when comparing an
    upgrade or checking the party's condition before leaving town.
  \item[Item:] Choose an owner and an object from that owner's inventory.
    The available actions include using, transferring, and discarding
    ordinary items. An item's available actions depend on the object.
  \item[Equip:] Select a companion and assign suitable weapons, armor,
    shields, and headgear. Check the resulting attack and defense values
    after making a change.
  \item[Tactics:] Manage party-related commands and battle behavior.
    Party arrangement determines the companions travelling at the front.
  \item[Door:] Open the door directly ahead when the party has the required
    access. Approach the door square and face it before using the command.
  \item[Search:] Examine the current position or the object directly ahead.
    This command is useful for chests, drawers, pots, and specific places
    on the ground.
\end{description}

A character-selection window identifies the companion whose spell,
equipment, or inventory is being used. Confirm the owner carefully when
several companions carry similar objects. Giving a Medical Herb to a
front-line companion places that herb among the items that companion can
use during a battle.

\subsection*{Working through menus}

Shops, inns, churches, and storage counters lead through a short sequence
of questions. Select the service, choose the companion or object involved,
and review the price or destination before accepting. A purchase may also
ask which companion should receive the item.

When the party has accumulated several spells or items, move through the
list to inspect the available choices. Returning to the previous window
lets you choose another owner or action. This is useful when distributing
new equipment or preparing several companions for a dungeon.

The opening data menu handles starting an adventure and managing its save
slots. Copying creates another copy of the selected adventure record.
Deleting removes the chosen record. Check the displayed slot and character
name carefully before confirming either operation.

\section{Data and status screens: attributes, experience, and equipment}

\subsection*{Checking a companion}

The Status command gives a compact view of a companion's development and
current condition. Begin with the name and level, then check HP and MP.
Their current values show the resources available for the next trip, and
their maximum values show the companion's capacity when fully recovered.

A character with ample maximum MP can still be nearly exhausted after a
long dungeon. Check the current value before relying on another healing
spell or planning a magical escape. The same habit helps identify a
companion who needs treatment before the next battle.

Experience records progress toward the next level. Level gains can improve
attributes and add spells or abilities. Check the Spell menu after a useful
level gain to see which new field abilities are available. A priest's
experience advice gives the amount still needed for the next level.

\subsection*{Attributes and equipment}

Strength, agility, base defense, intelligence, and luck describe the
companion's current development. Strength contributes to physical attack
power. Agility influences the order of actions in battle. Base defense
contributes to protection from physical attacks. Intelligence affects a
monster companion's command obedience.

Attack and defense show the current results of character development and
equipment. When considering a purchase, inspect the intended user's
existing weapon or armor first. After equipping the replacement, check
the new total and decide where the old piece will be most useful.

Special equipment properties deserve a separate look. Protection from
certain attacks or status effects can be valuable in the next dungeon.
Some weapons have unusual attack patterns, and some objects provide a
spell-like effect when used as an item. Choose equipment with the area's
enemies and the companion's role in mind.

\subsection*{Growth and the active party}

Monster species have individual level caps. A companion approaching its
cap may already have most of its available abilities. Consider its healing,
support, equipment, and durability when deciding how long to keep it in
the active party.

After a recruitment or party change, inspect the new companion's status,
spells, and equipment. Supply suitable gear and place useful consumables
in its inventory. This makes the new member ready to contribute on the
next journey.

The opening data menu also provides message-speed and sound settings.
Choose a comfortable text speed so that shop prices, service prompts,
and conversations remain easy to follow.

\section{Dialogue, searching, drawers, pots, chests, and hidden objects}

\subsection*{Talking through a town}

Speak to residents as you explore. Their conversations introduce nearby
places, local problems, unusual visitors, and changes in the town. A
shopkeeper may offer a useful destination, while another resident explains
why a building or route matters.

Conversations can change after an event or at another time of day. Return
to the people involved in the current objective after completing its
immediate task. The next instruction, reward, or permission to proceed
often comes through that follow-up conversation.

Face the person before using Talk. Across a counter, face the clerk from
the customer side. When a person has moved to another room or a building
has closed, check the time of day and look around the surrounding area.

\subsection*{Searching houses and shops}

Drawers, pots, barrels, and chests can hold useful objects. Approach and
face the container, then use Search. Check accessible rooms above or
behind shops as well as the obvious ground-floor rooms. A small find can
save a purchase or supply an item for the next trip.

The opening ship offers an immediate example. Search the drawer in the
cabin for a Medical Herb. In Santa Rosa, household searches provide early
supplies and equipment. These discoveries make it worth examining a new
room before leaving it.

Keep some inventory space available while exploring. Transfer spare
objects between companions or deposit them before a long treasure run.
A party carrying a full assortment of old equipment can run short of
space just as a useful new object is found.

\subsection*{Ground searches and unusual containers}

Some objects lie on a particular square of ground. Stand on the indicated
square and use Search. Details such as a flower patch, grave, well,
bucket, or conspicuous corner can help identify the correct position.
NPC dialogue may point toward the search area.

During a return visit, check rooms that have become accessible through
a new key or event. A newly opened cellar or side room may contain several
objects close together. Finish the searches in that room before returning
to the main route.

Some containers begin a fixed encounter. Keep the party reasonably healed
while opening unfamiliar containers in a dungeon, and have a recovery
item ready in an active companion's inventory.

\section{Day and night: access, shops, events, and changing time}

\subsection*{A town at two different times}

Day and night affect the daily life of towns. Residents move, businesses
open or close, and certain people offer their services at a particular
time. The change in lighting also makes it easy to recognize which part
of the day the party has reached.

Oracleberry provides a useful early example. Its special merchant conducts
business at night, and the fortune teller also has a nighttime routine.
Visit the town during both periods to see the available services and
conversations.

When a person is absent or a counter is closed, check nearby rooms and
consider another time of day. A resident may explain the schedule. Some
story events have their own timing and require a conversation or action
before the party can continue.

\subsection*{Letting time pass}

Travel around the countryside until the lighting changes, then enter
the town at the desired time. Choose a manageable area near the entrance
when waiting for night. Encounters during the walk give the party a chance
to earn experience and gold, so keep healing supplies available.

An ordinary stay at an inn restores the party and brings the day around
to morning. A story-directed stay may include an event. Once the party
is free to move again, check the surroundings and continue the current
objective.

\subsection*{Changing the time directly}

The Day-Night spell costs 3 MP and switches between day and night when
used in an appropriate field location. Select a companion who knows it
through the Spell command. This becomes convenient when returning to a
shop or searching for a person with a particular schedule.

The Lamp of Darkness changes day into night. Use it through the owner's
Item menu when that effect is needed. Enter or revisit the town afterward
to reach the nighttime activity.

\section{Doors, keys, moving obstacles, floor hazards, and puzzles}

\subsection*{Opening a route}

Approach a closed door and face it, then use Door. Locked doors require
the relevant key or learned unlocking ability. Keep important access
items available to the travelling party.

The Thieves' Key knowledge opens the ordinary locks associated with that
ability. The Reinhart Key is used for the castle's relevant locks. The
Magic Key opens its red-framed doors, and the Final Key provides access
to prison cells and other remaining key-locked doors.

A new key can make several older destinations worth revisiting. Complete
the newly accessible searches in one town before moving on to the next.
An entrance controlled by a story event or a mechanism requires that
particular action as well.

\subsection*{Rocks, switches, and connected passages}

Some movable obstacles can be pushed by walking into them from a suitable
side. Check the clear space around the object before moving it. A passage
can depend on where an obstacle ends up, so pay attention to nearby walls,
doors, and floor connections.

Inspect visible switches, levers, and unusual structures. Look for the
gate, track, staircase, or passage affected by the mechanism. After an
action, examine the changed area before moving farther away.

In dungeons with holes, sliding floors, or one-way connections, keep track
of the stairs and entrances that reconnect the rooms. A fall can lead to
a different section of the same dungeon. Locate the return path before
spending too many healing resources exploring its branches.

\subsection*{Crossing damaging terrain}

Hazardous floors can reduce HP while the party moves. Check the party's
condition after crossing them, especially when a battle may begin nearby.
Poison can also drain HP during travel and deserves prompt treatment.

Stepguard costs 2 MP and protects the party from damaging floor tiles.
Cast it before crossing an appropriate hazard. Check the party's condition
as the route moves into another area, and renew the protection when needed.

If the caster is short of MP, choose a route with fewer hazardous steps
and heal after the crossing. Keeping the party healthy before opening a
door or entering a new room gives more room to handle an unexpected battle.

\section{Inventory ownership, item transfer, capacity, and equipment use}

\subsection*{The twelve-item inventory}

Each companion has twelve inventory slots. Equipment occupies those
slots along with consumables, spare gear, and other carried objects.
Separate copies of a consumable occupy separate slots.

Leave room for discoveries before entering a treasure-rich area. Deposit
unused possessions, sell ordinary spare gear when it is safe to do so,
and share the carrying load among available companions. Keep the active
party's recovery supplies convenient to reach.

Item ownership matters in battle. A companion uses objects from its own
inventory during its turn. Give healing or status-recovery items to active
members who can use them when needed. A companion with a useful reusable
item should also be available for the encounter where it will be used.

\subsection*{Transferring and equipping}

Open Item, select the current owner, choose the object, and select the
transfer action. Choose the recipient and confirm the transfer. Check
the recipient's remaining space when moving several objects at once.

Use Equip to assign a suitable weapon, armor piece, shield, or headgear.
After a change, inspect the companion's attack and defense values. The
old equipment remains a possession until it is transferred, sold, stored,
or otherwise removed.

An older weapon or armor piece may suit another companion. Check the
recipient's equipment options before selling it. When moving equipped
gear, review both companions' equipment afterward so that each has the
intended loadout.

\subsection*{Using objects and conserving valuables}

Select Use from the object's available item actions. The game may ask
for a recipient or apply the effect to the field. Medical Herbs provide
healing, Antidotes treat poison, and Full Moon Herbs treat paralysis.
Their availability can help preserve a healer's MP during a long trip.

Certain weapons and other objects produce a spell-like effect when used.
Their useful action may be available alongside their equipment function.
Inspect the object's effect and give it to a companion who can make good
use of it.

Keep important keys, event objects, and unique possessions safe. Store
objects that are temporarily unnecessary when the depository accepts
them. Resolve cursed equipment through the appropriate church service
before trying to change an affected loadout.

\section{Inns, churches, saving, revival, and status recovery}

\subsection*{Resting before the next journey}

An inn restores the living party's HP and MP. Speak to the innkeeper,
review the cost, and accept the stay. Free resting opportunities are
particularly useful during early training or while gathering enough gold
for an equipment upgrade.

Check the party's condition after resting. Treat persistent ailments and
revive fallen companions before setting out. A fully supplied healer and
healthy front line make the next journey much more comfortable.

\subsection*{Church services}

Speak to the priest and select the service needed:
\begin{description}
  \item[Record Adventure:] Save the current adventure. Confirm the chosen
    record and wait for the priest to finish the saving dialogue.
  \item[Experience advice:] Learn how much experience a companion needs
    to reach the next level.
  \item[Revival:] Restore a fallen companion to life. Select the companion
    and review the required payment.
  \item[Poison treatment:] Remove poison from the selected companion.
  \item[Curse removal:] Resolve a curse affecting a companion or its
    equipment through the priest's service.
\end{description}

Save after an important acquisition, a difficult dungeon, or a substantial
party change. Before a longer expedition, visit the inn and church so
that the saved record reflects the prepared party.

\subsection*{Recovery away from town}

Healing spells and consumables let the party recover between battles.
Antidote costs 2 MP and treats poison. Numboff costs 2 MP and treats
paralysis. Carrying Antidotes and Full Moon Herbs gives additional options
when the appropriate caster is unavailable or short of MP.

Vivify costs 10 MP and attempts to revive a fallen ally. A failed attempt
requires another cast, so allow enough MP when relying on it. Revive costs
20 MP and restores a fallen ally with full HP. A Leaf of the World Tree
also provides revival and is a valuable emergency possession.

After a revival, check HP, MP, and remaining supplies. A recovery spell
can make it possible to finish a short route or reach the exit. A return
to town provides a comfortable chance to rest, replenish items, and save.

\subsection*{Recovering from defeat}

When the party is defeated, half of the gold being carried is lost.
Banked gold remains protected. Check the party's condition after returning
to the recovery point, then arrange revival, treatment, and rest.

Revisit the equipment and supply plan before trying the area again.
A protective item, a different active companion, or an extra recovery
object may give the party the help it needs on the next attempt.

\section{Banks, gold storage, shopping, and selling}

\subsection*{Depositing gold and possessions}

The depository offers storage for gold and items. Oracleberry has one
near Monster Gramps. Gold deposits are made in increments of 1,000 G.
Keep enough spending money for the next inn stay, treatment, and planned
purchases, then deposit the amount that can be set aside.

Stored gold is protected from the loss caused by defeat. Withdraw it when
an equipment purchase or recovery expense is due. Keeping a reserve in
the bank makes rebuilding a defeated party easier.

The item-storage service is useful for spare equipment and possessions
that are being kept for later. Select the owner and object when depositing,
then use the withdrawal service when the object is needed again. Leave
space with the intended recipient before collecting several stored items.

When a monster is left with Monster Gramps, its carried equipment and
items are transferred to the depository. Retrieve the needed possessions
and equip the companion again when bringing it back into the travelling
party.

\subsection*{Making a purchase}

Speak to the shopkeeper, choose the item, and review its price and the
intended recipient. Consider which active companion gains the most from
the upgrade. A dependable weapon for a frequent attacker or protective
gear for a vulnerable member can have an immediate effect on the next trip.

After buying equipment, assign it through Equip and check the resulting
status values. Transfer useful older gear to another eligible companion.
Keep enough cash for recovery and essential consumables after the purchase.

Oracleberry's nighttime merchant sells the wagon for 300 G. Having the
required cash ready makes that visit convenient. Purchases tied to travel
or party access are worth accounting for alongside ordinary equipment.

\subsection*{Selling and keeping useful objects}

The selling service quotes a value before the transaction is completed.
Review the object and its owner carefully. Ordinary gear that has been
replaced throughout the party can provide money and free inventory space.

Keep unique objects, useful reusable items, and equipment intended for a
companion who will be trained later. The depository provides space for
possessions that are worth retaining. Store them together with a clear
idea of which companion or journey they will support.

\section{Travel progression: walking, wagon, ship, carpet, castle, and dragon}

\subsection*{Walking and the wagon}

On foot, follow paths through the countryside and look for the entrances
to towns, caves, towers, and other destinations. Mountains and other
obstacles define the available routes. Encounters along the way consume
healing supplies and MP, so consider the journey to the entrance when
preparing for a dungeon.

The wagon follows the party during appropriate field travel and provides
room for additional companions. Arrange the active members and distribute
their equipment before entering an area with restricted wagon access.
Keep the supplies needed inside that area with the companions who will
be available there.

\subsection*{Sailing}

The ship opens sea routes between continents and islands. Board at its
docking point and sail toward the desired coast. Choose an accessible
landing place near the destination, then continue on foot.

Sea travel has its own encounters. Maintain enough healing for the voyage
and the exploration after landing. A nearby town or inn can provide a
convenient starting point for a longer island trip.

Remember the landing point before entering a dungeon or making a detour.
The coastline and nearby landmarks help locate the ship again when the
party is ready to continue.

\subsection*{Using the Magic Carpet}

Use the Magic Carpet from its owner's inventory on suitable open ground.
It carries the party across plains and water, making several routes more
direct. Forests and mountains constrain its movement.

Choose a clear approach to the destination and a suitable place to land.
After dismounting, walk into the town or dungeon. While travelling by
carpet, X opens the world map and helps identify nearby land and routes.

\subsection*{The flying castle}

Once Zenithia is available for travel, its movement opens another way to
cross the world. Use the directional pad while controlling the castle.
Press X to view the world map and check the destination before choosing
the approach.

The castle's interior also provides a place to regroup during that part
of the adventure. Keep the party ready for the destination reached after
leaving it, especially when continuing directly into a dungeon.

\subsection*{Dragon travel}

The Zenithian Bells summon the flying ride when they are available.
Use the item outdoors and choose the desired route through the world.
This form of flight provides access across terrain such as forests and
mountains.

Choose a suitable landing point near the destination, then dismount and
continue on foot. Check the active party and healing supplies before
entering the next area. A fast flight can place the party close to a
demanding dungeon with little travel in between.

\section{Return destinations, Outside, maps, and safe backtracking}

\subsection*{Return and Outside}

Return costs 8 MP and transports the party to an available destination
in its list. Towns and other eligible stops become available through the
journey. Select the desired destination from the Spell command while
standing in a suitable outdoor location.

Outside costs 8 MP and takes the party out of a dungeon that allows
magical escape. It is useful after collecting treasure or when healing
resources are running low. Once outside, the party can choose the next
overland route or use Return.

If the same companion will cast both spells, keep at least 16 MP available
for that trip home. Allow additional MP for any healing still needed before
reaching the exit. The remaining value on the Status screen is useful
when deciding whether to explore another branch.

\subsection*{Other travel aids}

A Wing of Wyvern provides an item-based journey back to the last town
visited. Carry one with an available companion and use it from a suitable
outdoor position when that destination is useful.

Repel costs 4 MP and helps avoid random battles with weak enemies. Holy
Water also helps ward off weak encounters. These aids can make a familiar
journey more convenient as the party grows stronger.

Use the Wayfarer's Map to examine the overworld when the item is available.
Compare the party's position with the coast, rivers, mountains, and nearby
settlements. The map is particularly useful when a new vehicle opens a
route across water or around a large obstacle.

\subsection*{Returning for earlier discoveries}

A new key or travel ability can open several useful return visits. Choose
one town or region, collect its newly accessible treasure, and then move
to the next nearby stop. Visit shops that have become useful to the
current party and retrieve any stored equipment needed for the trip.

Keep a short mental route through the area: the landing or Return point,
the rooms to search, the services to use, and the way back to the vehicle
or next destination. This makes a return visit comfortable even when the
town was last explored much earlier in the adventure.

Finish a longer outing by treating ailments, restoring HP and MP,
reorganizing useful possessions, and recording the adventure at a church.
The party is then ready to begin its next journey from a convenient stop.

